\section{Discussion}\label{sec:discussion}

\subsection{What the results do and do not say}

The results say that the within-CDE rural penalty is precisely zero at
the median, null at both the extensive and intensive margins, and small but
imprecisely estimated in the mean; and that the aggregate rural
penalty---which is real---is located in between-CDE composition. They do
not say that rural credit markets are no different from urban ones; they
say the NMTC's deployment data do not detect such a difference once the
intermediary is held fixed. They also do not fully adjudicate \emph{why}
composition looks as it does: reading the between-CDE component as
intermediary selection is the natural interpretation, but rural market
conditions could in principle drive low-leverage intermediaries to
specialize in rural deployment, which would put part of the between-CDE
component back on the market-structure side of the ledger. Distinguishing
those readings requires variation in CDE--market assignment that the
public release does not contain. Nor, finally, do the results say that any
particular CDE could not mobilize more in rural tracts; the estimates are
averages over the fixed-effects distribution.

\subsection{Policy reading}

Reading the evidence this way carries three implications, in ascending
order of specificity. First, the marginal rural deal does not appear
constrained by market structure under the existing program design: when a
given intermediary goes rural, its mobilization does not measurably
suffer. Second, the aggregate rural gap is therefore amenable to
\emph{intermediary-allocation} levers---shifting allocation toward CDEs
with demonstrated high-leverage capacity, independent of their current
rural orientation---without redesigning the credit itself. Third, the 20\%
mandate as currently observed operates on a margin (cumulative deployment
share) where behavior shows no bunching; if the policy goal is rural
mobilization rather than rural deal counts, the mandate is pointed at the
wrong margin.

\subsection{Research program}

This paper is the first phase of a sequence: the NMTC alone here; a
multi-program U.S. comparison (NMTC, LIHTC, Opportunity Zones) next; and an
international extension to development-bank project finance where the same
mobilization framework applies. The LIC-eligibility regression
discontinuity, which requires the ACS tract merge, is the immediate
companion piece.
